\section{Introduction}\label{sec:intro}

Does there exist a $3\times3$ magic square whose nine entries are distinct
positive integer squares? LaBar posed this question in 1984; Gardner
popularised it and offered \$100 for an example (see the historical account
in Bremner~\cite{BRE99}, \S0). Two classical near misses approach the
question in different ways: Bremner's fully magic array has seven square
entries~\cite{BRE99}, whereas Sallows's nine-square array has seven equal
ordinary line sums~\cite{SAL97}. Here an ordinary line means a row, column,
or long diagonal. We call an array an \emph{$(8,7)$ array} if its nine
entries are positive and pairwise distinct, exactly eight are integer
squares, and exactly seven ordinary lines have one common sum. A
\emph{$(9,7)$ array} satisfies the same conditions with all nine entries
square.

The common structure is elementary but arithmetically useful: an array
with seven equal line sums consists of three arithmetic progressions
with one common difference. Closing its eighth line asks that their three
centres form another arithmetic progression. This separates two questions:
how to construct near misses, and which differences can support a fully
magic interaction. The prime support in this paper belongs to the
\emph{difference}, not to the nine entries.

\subsection{Contributions and their dependencies}

\paragraph{Complete fixed-support results.}
Our most concrete classification is Theorem~\ref{main:smoothsupport}.
Using the complete $S$-unit totals of von K\"anel--Matschke~\cite{VKM16},
we give a bidirectional reconstruction of the relevant primitive right
triangles and square progressions, test every possible centre interaction,
and classify the resulting $(9,7)$ arrays. For primes at most $19$ this
gives no fully magic interaction, exactly $384$ classes modulo the
dihedral group and common rational-square scaling, and the unique minimum
\[
\eta=\frac{|\text{common line sum}-\text{failed line sum}|}
{|\text{transversal difference}|}
=\frac{347984603}{9837828000}.
\]
The imported total certifies completeness beyond the finite production
bound. The reconstruction, classification and optimization are the finite
deductions proved here. In particular, the small-support example already
displayed by Houston~\cite{HOU16B} is credited as such; the atlas
classifies the entire stated domain. Separately, the seven triangles in
Aebi's classification~\cite{AEB26} have distinct area squareclasses. We
deduce that a difference with at most three prime factors supports at most
one positive square progression (Theorem~\ref{main:threeprime}).

\paragraph{Constructive results.}
We determine the integral image lattice of the line-sum map and prove its
Smith form, refining the classical normal form
(Proposition~\ref{prop:linesumsnf}). Starting with catalogued
Bremner--Wesolowski carriers, we then prove a uniform retraction formula,
classify every pair-closing shift and determine the least one. A second
infinite construction lifts an elliptic orbit to pairwise inequivalent
primitive $(8,7)$ arrays; applying Cowan's equidistribution
theorem~\cite{COW20} gives the density of the admissible \emph{orbit
indices}. A finite classification of the difference-$840$ bridge supplies
a further explicit construction. These results form
Theorem~\ref{main:families}: two infinite mechanisms and one finite bridge,
with the classical carriers and elliptic ingredients attributed at use.

\paragraph{The interaction and its arithmetic restrictions.}
We develop the centre-interaction condition in two settings. For fixed
difference $840$ it is a relation on three doubled elliptic points.
Allowing the difference to vary and dividing out the common root scale
gives a different, explicitly related surface. We prove that this
normalized surface is irreducible and that its degree-five reconstruction
over the square-progression cone has generic Galois group $S_5$
(Theorem~\ref{main:surface}). For a prescribed
difference an additional rational-square condition on the scale is
essential. We specialize Caro--Garc\'ia-Fritz's theorem~\cite{CGF25} to
obtain non-effective fixed-squareclass finiteness, close the centre-$841$
fibre and a fixed Bremner shadow by explicit elliptic quotients, and give
a countercertificate to deciding the ambient interaction from the three
separate finite Kummer classes alone. Finally we derive a nondegenerate
five-term $S$-unit equation, local valuation restrictions and precise
reductions of the remaining four-prime branches. The finite-support
catalogue above is a separate effective argument; it does not depend on
the non-effective finiteness bound.

\subsection{Relation to earlier work and organization}

The normal form and its interpretation by square progressions occur in
Houston~\cite{HOU16A,HOU16B}, following Lucas~\cite{LUC76,LUC91}; the
equal-area triangle correspondence is classical and is used by
Robertson~\cite{ROB96}. Bremner~\cite{BRE99,BRE01}, Boyer's
survey and catalogue~\cite{BOY05,BOYWEB}, and Fituvalu's
tables~\cite{FIT} supply the carrier context. The support law uses the
unit-equation theorem of Evertse--Schlickewei--Schmidt~\cite{ESS02}.
These inputs are distinguished from the refinements, constructions and
applications proved in this paper. Results restricting the prime factors
of the \emph{entries}, due to Rabern, Woll, Labruna and
Weisenberg~\cite{RAB03,WOL18,LAB18,WEI23}, concern a complementary
quantity; Appendix~\ref{app:literature} records their scope and further
bibliographic comparisons. It also explains the elementary quartic-field
completion of Bremner's example, which shows why a shape-only obstruction
cannot settle the rational problem. The historical label ``Parker'' in
some companion records recalls Parker's popular near miss~\cite{PARKER};
it neither attributes LaBar's problem to him nor identifies his particular
array with the configurations studied here.

Section~\ref{sec:normalform} fixes the normal form and conventions.
Section~\ref{sec:smoothsupport} proves the support exclusions and exact
atlas, followed by the constructions in Section~\ref{sec:families}.
Sections~\ref{sec:surface} and~\ref{sec:finiteness} treat the interaction
geometry and specific arithmetic obstructions.
Sections~\ref{sec:supportlaw} and~\ref{sec:frontier} develop the local
support laws and four-prime frontier; Section~\ref{sec:program} identifies
the remaining questions. The appendices retain the bounded census,
least-shift details, supplementary reductions, source comparisons and
reproducibility information. External rank and point-list computations
are identified where they enter a proof; their archived output records
are distinct from the exact computations rerun by the public verifiers.
